\documentclass[10pt, letterpaper]{article}

% Packages:
\usepackage[
    ignoreheadfoot, % set margins without considering header and footer
    top=2 cm, % seperation between body and page edge from the top
    bottom=2 cm, % seperation between body and page edge from the bottom
    left=2 cm, % seperation between body and page edge from the left
    right=2 cm, % seperation between body and page edge from the right
    footskip=1.0 cm, % seperation between body and footer
    % showframe % for debugging 
]{geometry} % for adjusting page geometry
\usepackage{titlesec} % for customizing section titles
\usepackage{tabularx} % for making tables with fixed width columns
\usepackage{array} % tabularx requires this
\usepackage[dvipsnames]{xcolor} % for coloring text
\definecolor{primaryColor}{RGB}{0, 79, 144} % define primary color
\usepackage{enumitem} % for customizing lists
\usepackage{fontawesome5} % for using icons
\usepackage{amsmath} % for math
\usepackage[
    pdftitle={Neo Cass's CV},
    pdfauthor={Neo Cass},
    colorlinks=true,
    urlcolor=primaryColor
]{hyperref} % for links, metadata and bookmarks
\usepackage[pscoord]{eso-pic} % for floating text on the page
\usepackage{calc} % for calculating lengths
\usepackage{bookmark} % for bookmarks
\usepackage{lastpage} % for getting the total number of pages
\usepackage{changepage} % for one column entries (adjustwidth environment)
\usepackage{paracol} % for two and three column entries
\usepackage{ifthen} % for conditional statements
\usepackage{needspace} % for avoiding page brake right after the section title
\usepackage{iftex} % check if engine is pdflatex, xetex or luatex

% Ensure that generate pdf is machine readable/ATS parsable:
\ifPDFTeX
    \input{glyphtounicode}
    \pdfgentounicode=1
    % \usepackage[T1]{fontenc} % this breaks sb2nov
    \usepackage[utf8]{inputenc}
    \usepackage{lmodern}
\fi



% Some settings:
\AtBeginEnvironment{adjustwidth}{\partopsep0pt} % remove space before adjustwidth environment
\pagestyle{empty} % no header or footer
\setcounter{secnumdepth}{0} % no section numbering
\setlength{\parindent}{0pt} % no indentation
\setlength{\topskip}{0pt} % no top skip
\setlength{\columnsep}{0cm} % set column seperation
\makeatletter
\let\ps@customFooterStyle\ps@plain % Copy the plain style to customFooterStyle
\patchcmd{\ps@customFooterStyle}{\thepage}{
    \color{gray}\textit{\small Neo Cass - Page \thepage{} of \pageref*{LastPage}}
}{}{} % replace number by desired string
\makeatother
% \pagestyle{customFooterStyle}

\titleformat{\section}{\needspace{4\baselineskip}\bfseries\large}{}{0pt}{}[\vspace{1pt}\titlerule]

\titlespacing{\section}{
    % left space:
    -1pt
}{
    % top space:
    0.3 cm
}{
    % bottom space:
    0.2 cm
} % section title spacing

\renewcommand\labelitemi{$\circ$} % custom bullet points
\newenvironment{highlights}{
    \begin{itemize}[
        topsep=0.10 cm,
        parsep=0.10 cm,
        partopsep=0pt,
        itemsep=0pt,
        leftmargin=0.4 cm + 10pt
    ]
}{
    \end{itemize}
} % new environment for highlights

\newenvironment{highlightsforbulletentries}{
    \begin{itemize}[
        topsep=0.10 cm,
        parsep=0.10 cm,
        partopsep=0pt,
        itemsep=0pt,
        leftmargin=10pt
    ]
}{
    \end{itemize}
} % new environment for highlights for bullet entries


\newenvironment{onecolentry}{
    \begin{adjustwidth}{
        0.2 cm + 0.00001 cm
    }{
        0.2 cm + 0.00001 cm
    }
}{
    \end{adjustwidth}
} % new environment for one column entries

\newenvironment{twocolentry}[2][]{
    \onecolentry
    \def\secondColumn{#2}
    \setcolumnwidth{\fill, 4.5 cm}
    \begin{paracol}{2}
}{
    \switchcolumn \raggedleft \secondColumn
    \end{paracol}
    \endonecolentry
} % new environment for two column entries

\newenvironment{header}{
    \setlength{\topsep}{0pt}\par\kern\topsep\centering\linespread{1.5}
}{
    \par\kern\topsep
} % new environment for the header

\newcommand{\placelastupdatedtext}{% \placetextbox{<horizontal pos>}{<vertical pos>}{<stuff>}
  \AddToShipoutPictureFG*{% Add <stuff> to current page foreground
    \put(
        \LenToUnit{\paperwidth-2 cm-0.2 cm+0.05cm},
        \LenToUnit{\paperheight-1.0 cm}
    ){\vtop{{\null}\makebox[0pt][c]{
        \small\color{gray}\textit{Last updated in June 2024}\hspace{\widthof{Last updated in June 2024}}
    }}}%
  }%
}%

% save the original href command in a new command:
\let\hrefWithoutArrow\href

% new command for external links:
\renewcommand{\href}[2]{\hrefWithoutArrow{#1}{\mbox{\ifthenelse{\equal{#2}{}}{ }{#2 }\raisebox{.15ex}{\footnotesize \faExternalLink*}}}}


\begin{document}
    \begin{header}
        \textbf{\fontsize{22 pt}{22 pt}\selectfont Neo Michael Cass}

        \vspace{0.3 cm}

        \normalsize
        \mbox{{\color{black}{\footnotesize\faMapMarker*}\hspace*{0.13cm}Geneva, CH}}
        \kern 0.5 cm
        \mbox{\hrefWithoutArrow{https://linkedin.com/in/neocass}{\color{black}{\footnotesize\faLinkedinIn}\hspace*{0.13cm}neocass}}
        \kern 0.5 cm
        \mbox{\hrefWithoutArrow{https://gitlab.cern.ch/ncass}{\color{black}{\footnotesize\faGitlab}\hspace*{0.13cm}ncass}}
        \kern 0.5 cm
        % \mbox{\hrefWithoutArrow{https://github.com/neo-cass}{\color{black}{\footnotesize\faGithub}\hspace*{0.13cm}neo-cass}}
        % \kern 0.5 cm
    \end{header}

    \vspace{- 0.0 cm}
    % \vspace{0.001 cm}

    \section{Experience}


        \begin{twocolentry}{
        \textit{Meyrin, Switzerland}    
            
        \textit{Sep. 2024 to Now}}
            \textbf{Software Engineer, Origin Graduate}
            
            \textit{CERN | BE-CEM-MTA}
        \end{twocolentry}

        \vspace{0.10 cm-3px}
        
        \begin{onecolentry}
            \begin{highlights}
                \item Dockerised the compilation and development environment of the CMW Wrapper, a shared library which allows CMW to be used from LabVIEW, for several platforms. 
                \item Moved the compilation and deployment of the libraries to GitLab CI/CD, greatly reducing the time to create a release.
                \item Fixed numerous errors in the CMW Wrapper and ran stress tests invoked by Python to ensure the new release is stable for clients. 
                \item Continued maintenance of the legacy RADE JAPC Service and commissioning of the new service. 
                \item Presented my Master's Thesis and the current state of the RADE JAPC Service to BE-CEM. 
                \item Attended the CERN School of Computing, improving my networking skills and core Computer Science competencies.
               
            \end{highlights}
        \end{onecolentry}


        \vspace{0.2 cm-3px}
        
        \begin{twocolentry}{
        \textit{Bath, UK}    
            
        \textit{Oct. 2023 to May 2024}}
            \textbf{Software Engineer, Student External User}
            
            \textit{CERN | University of Bath}
        \end{twocolentry}

        \vspace{0.10 cm-3px}
        \begin{onecolentry}
            \begin{highlights}
                \item Continued my work on the RADE JAPC Service as part of my Master's Thesis. 
                \item Significantly improved the interface by creating a REST API for JAPC get/set operations and implemented an analogous WebSocket interface for JAPC subscriptions. 
                \item Benchmarked the service to evaluate the cost of all the architectural changes. The new service outperformed or matched the legacy service.
                \item Implemented test clients in Python and LabVIEW and saw that these have an average of 85\% fewer lines of code than the legacy service, therefore reducing the dependency footprint on end-users and increasing maintainability.
            \end{highlights}
        \end{onecolentry}


        \vspace{0.2 cm-3px}

        \begin{twocolentry}{
        \textit{Meyrin, Switzerland}    
            
        \textit{Aug. 2022 to Sept. 2023}}
            \textbf{Software Engineer, Technical Student}
            
            \textit{CERN | BE-CEM-MTA}
        \end{twocolentry}

        \vspace{0.10 cm-3px}
        \begin{onecolentry}
            \begin{highlights}
                \item Rewrote the legacy RADE JAPC Service, which allows devices to be monitored and controlled with JAPC from LabVIEW.
                \item New service is a Dockerised Spring Boot application, built with Gradle on GitLab CI/CD and deployed on Openshift. The new system is much more modern, scalable and maintainable. Notably, the new service has 75\% less lines of code and can support additional client languages, such as Python. 
                \item Responsible for maintaining and upgrading the legacy service in parallel, which is used in portable test benches (SM18, FRESCA), AD and expert tools from QPS, FGC, PIC, MPP, cryo and kickers. 
            \end{highlights}
        \end{onecolentry}



    
    \section{Technologies}



        
        \begin{onecolentry}
            \textbf{Languages:} Java, Python, C++, C, Haskell, SQL and LabVIEW. \textbf{More importantly}: I understand the fundamental principles that underpin programming so I can quickly respond to the challenge of new languages. 
        \end{onecolentry}

        \vspace{0.2 cm}

        \begin{onecolentry}
            \textbf{Software:} Git and GitLab CI/CD, Spring Framework and Spring Boot, Apache Tomcat, Docker, OpenShift and Gradle. \textbf{Experienced developing services with:} JAPC, CMW and RBAC.
        \end{onecolentry}


    
    \section{Education}



        
        \begin{twocolentry}{
            
            
        \textit{Sept. 2019 to June 2024}}
            \textbf{University of Bath}

            \textit{MComp(Hons) in Computer Science and Mathematics}
        \end{twocolentry}

        \vspace{0.10 cm-3px}
        \begin{onecolentry}
            \begin{highlights}
                \item \textbf{Overall Grade}: First-Class Honours (71\%) | \textbf{Master's Grade}: First-Class Honours (80\%)
                \item Master's Thesis: \textit{RADE JAPC Service: A New Service For Controlling and Monitoring Devices at CERN}
                \item Bachelor's Thesis: \textit{Detecting Phishing Emails with Ensemble Learning}
                \item Programme Summary: Problem solving and team-working skills are developed throughout, while gaining a profound understanding in fundamental Computer Science and Mathematics. 
            \end{highlights}
        \end{onecolentry}


        % \vspace{0.2 cm-3px}

        % \begin{twocolentry}{
            
            
        % \textit{Sept. 2017 to June 2019}}
        %     \textbf{Callywith College}

        %     \textit{A-Levels}
        % \end{twocolentry}

        % \vspace{0.10 cm-3px}
        % \begin{onecolentry}
        %     \begin{highlights}
        %         \item Maths (A*) Computer Science (A) Physics (A) | AS-Levels: Further Maths (A) English Language (B)
        %     \end{highlights}
        % \end{onecolentry}


        % \vspace{0.2 cm-3px}

        % \begin{twocolentry}{
            
            
        % \textit{Sept. 2012 to June 2017}}
        %     \textbf{Looe Community Academy}

        %     \textit{GCSEs}
        % \end{twocolentry}

        % \vspace{0.10 cm-3px}
        % \begin{onecolentry}
        %     \begin{highlights}
        %         \item 8 (A-A*) including: English Literature/Language (A*/A) Mathematics (A*)
        %     \end{highlights}
        % \end{onecolentry}



    

\end{document}